%! Logarithmic Function Manipulation — Unit 5, Lesson 5.4 — Warm-Up (Answer Key)
\documentclass[10pt]{article}
\usepackage{precalculus-article}
\usepackage{precalculus-key}   % pulls in -boxes; do NOT also load -boxes
\newcommand{\TallMath}[1]{$\displaystyle #1\rule[-1.4em]{0pt}{3.2em}$}

\begin{document}

\pageheader{Unit 5, Lesson 5.4}{Warm-Up --- Answer Key}
\namedateperiod

\vspace{0.15in}

\begin{enumerate}[leftmargin=1.8em, itemsep=14pt]

\item Evaluate each logarithm.

\vspace{6pt}
\begin{center}
\renewcommand{\arraystretch}{2.0}
\begin{tabular}{@{} c @{\qquad} c @{\qquad} c @{}}
  $\log_2(32) = $ \ans{5} &
  $\log_3(81) = $ \ans{4} &
  $\log_{10}(1000) = $ \ans{3} \\
\end{tabular}
\end{center}

\item Let $f(x) = \log_5(x)$. Evaluate $f(25)$, $f(5)$, and $f(1)$.

\vspace{6pt}
\begin{center}
\renewcommand{\arraystretch}{2.0}
\begin{tabular}{@{} c @{\qquad} c @{\qquad} c @{}}
  $f(25) = $ \ans{2} &
  $f(5) = $ \ans{1} &
  $f(1) = $ \ans{0} \\
\end{tabular}
\end{center}

\vspace{4pt}
What pattern do you notice in the outputs as the input multiplies by 5?

\ansline{The output increases by 1 each time the input multiplies by 5 --- additive change.}

\item Recall the exponent rule: $2^3 \cdot 2^4 = 2^{\ans{7}}$.

\vspace{6pt}
Based on this, predict: $\log_2(8 \cdot 16) = \log_2(8) + \log_2(16) = $ \ans{3} $+$ \ans{4} $=$ \ans{7}.

\vspace{4pt}
Verify by evaluating $\log_2(128)$ directly: $\log_2(128) = $ \ans{7}.

\end{enumerate}

\end{document}
